\chapter{A hash commitment instead of the ideal one}
\label{chap:commitment}

In this chapter the commitment functionality of Chapter~\ref{chap:jko} is replaced by the hash
commitment $c = G(Z, r)$, where $G \colon V \times R \to D$ is a random oracle and $r$ is uniform in
a finite set $R$. The oblivious transfer stays ideal. Zero knowledge becomes statistical, with a
simulator that programs the oracle, and soundness becomes extraction plus a binding error.

\section{Oracle runs}

\begin{definition}[Runs]
  \label{def:ans}
  \lean{GcZk.Com.ans, GcZk.Com.Hits}
  \leanok
  A query strategy chooses each query from the previous answers. A run queries a point $p$ if one
  of its queries is $p$.
\end{definition}

\begin{lemma}[Reprogramming an unqueried point]
  \label{lem:ans-update}
  \lean{GcZk.Com.ans_congr, GcZk.Com.ans_update}
  \leanok
  \uses{def:ans}
  If a run does not query $p$, reprogramming the oracle at $p$ changes none of its answers.
\end{lemma}

\begin{proof}
  \leanok
  Induction on the number of queries.
\end{proof}

\section{Zero knowledge}

\begin{definition}[Malicious verifier and the three games]
  \label{def:com-games}
  \lean{GcZk.Com.MalV, GcZk.Com.msg, GcZk.Com.zw, GcZk.Com.realFin, GcZk.Com.simFin,
    GcZk.Com.play, GcZk.Com.point, GcZk.Com.real, GcZk.Com.hyb, GcZk.Com.sim}
  \leanok
  \uses{def:view, lem:ans-update}
  A malicious verifier runs in three phases with adaptive queries to $G$: $q_1$ queries and then
  the material and the OT pairs; after the commitment, $q_2$ queries and then the opened seed;
  after the prover's last message, $q_3$ queries and then its output bit. In the \emph{real}
  execution the prover commits with $G(Z_w, r)$, where $Z_w$ is its output label, and sends
  $(Z_w, r)$ if the opening passes the check. The \emph{simulator} commits to a uniform $c$, and if
  the opening passes the check it computes $Z = L_{\mathrm{out}} + \Delta$ from the seed, sends
  $(Z, r)$ and programs $G(Z, r) := c$. It does not use the witness. The \emph{hybrid} uses the real
  commitment point $p = (Z_w, r)$, a uniform commitment $c$, and the oracle $G[p := c]$ after the
  commitment.
\end{definition}

\begin{lemma}
  \label{lem:sim-fin}
  \lean{GcZk.Com.simFin_eq}
  \leanok
  \uses{def:com-games, lem:prover-label}
  If $C(w) = 1$, the simulator's last message is the honest prover's. In the formalization the
  hypothesis is that the prover's check is sound (Lemma~\ref{lem:check-sound-ot}), which for the
  check of Definition~\ref{def:view} follows from $C(w) = 1$; the same hypothesis is used in
  Lemma~\ref{lem:hyb-sim}.
\end{lemma}

\begin{proof}
  \leanok
  Lemma~\ref{lem:prover-label}.
\end{proof}

\begin{lemma}[Real and hybrid]
  \label{lem:hyb-real}
  \lean{GcZk.Com.Bad₁, GcZk.Com.swap, GcZk.Com.ans₁_swap, GcZk.Com.point_swap,
    GcZk.Com.bad₁_swap, GcZk.Com.swap_swap, GcZk.Com.hyb_eq_real,
    GcZk.Com.card_hyb_eq_card_real}
  \leanok
  \uses{def:com-games}
  If the first phase does not query $p$, swapping the value $G(p)$ with $c$ turns the hybrid into
  the real execution. The swap is an involution on these pairs $(G, c)$, so the hybrid and the real
  execution output $1$ equally often on them.
\end{lemma}

\begin{proof}
  \leanok
  By Lemma~\ref{lem:ans-update} the first phase, hence $p$, is the same for $G$ and for $G[p := c]$;
  after the swap the real commitment is $G[p := c](p) = c$.
\end{proof}

\begin{lemma}[Hybrid and simulator]
  \label{lem:hyb-sim}
  \lean{GcZk.Com.Bad₂₃, GcZk.Com.hyb_eq_sim}
  \leanok
  \uses{def:com-games, lem:sim-fin}
  The hybrid and the simulator agree unless the second phase, or the third phase after an abort,
  queries $p$.
\end{lemma}

\begin{proof}
  \leanok
  After a successful opening both use $G[p := c]$ in the third phase, by Lemma~\ref{lem:sim-fin};
  elsewhere apply Lemma~\ref{lem:ans-update}.
\end{proof}

\begin{theorem}[Statistical zero knowledge]
  \label{thm:zk-rom}
  \lean{GcZk.Com.card_real_le, GcZk.Com.card_sim_le, GcZk.Com.card_bad₁, GcZk.Com.card_bad₂₃,
    GcZk.Com.sum_card_swap, GcZk.Com.sum_bad_le, GcZk.Com.zk_rom}
  \leanok
  \uses{lem:hyb-real, lem:hyb-sim}
  For every malicious verifier and every witness $w$ with $C(w) = 1$, the probabilities that the
  verifier outputs $1$ in the real execution and with the simulator differ by at most
  \[ \frac{q_1 + q_2 + q_3}{|R|} . \]
\end{theorem}

\begin{proof}
  \leanok
  By Lemmas~\ref{lem:hyb-real} and \ref{lem:hyb-sim}, the two probabilities differ by at most the
  probability of the bad events. None of the queries involved depends on $r$: before the opening
  the simulator does not use $r$, and after an abort it never reveals it. So each query hits
  $p = (Z_w, r)$ for at most one value of $r$.
\end{proof}

\section{Soundness}

\begin{definition}[Cheating prover and extractor]
  \label{def:bind}
  \lean{GcZk.Bind.Prover, GcZk.Bind.a₁, GcZk.Bind.qry₁, GcZk.Bind.com, GcZk.Bind.qry₂,
    GcZk.Bind.opn, GcZk.Bind.Accept, GcZk.Bind.ExtWins}
  \leanok
  \uses{def:ans}
  A cheating prover sees a view of the garbling, makes $q_1$ queries to $G$ and commits to $c$;
  after the opening it makes $q_2$ queries and opens a point. The verifier accepts if the opened
  point has the target label and its value is $c$. The extractor returns the label of the first
  first-phase query answered with $c$. It does not see the opening, so it is an adversary for the
  authenticity game.
\end{definition}

\begin{lemma}[Fiber bound]
  \label{lem:fresh}
  \lean{GcZk.Bind.card_fresh_le, GcZk.Bind.ans_update_of_avoid, GcZk.Bind.qry₁_update,
    GcZk.Bind.a₁_update, GcZk.Bind.qry₂_update, GcZk.Bind.opn_update}
  \leanok
  \uses{lem:ans-update}
  Let $E$ be a set of oracles, and let $p(G)$ be a point and $f(G)$ a value that, for $G \in E$, do
  not change when $G$ is reprogrammed at $p(G)$. Then $G \in E$ and $G(p(G)) = f(G)$ for at most a
  fraction $1/|D|$ of the oracles.
\end{lemma}

\begin{proof}
  \leanok
  The map $(G, x) \mapsto G[p(G) := x]$ is injective on the pairs with $G \in E$ and
  $G(p(G)) = f(G)$: the image determines $p(G)$, hence $G$ off $p(G)$, and $G(p(G)) = f(G)$ is
  determined too.
\end{proof}

\begin{lemma}
  \label{lem:bad-of-accept}
  \lean{GcZk.Bind.Coll, GcZk.Bind.Hit₂, GcZk.Bind.HitOpen, GcZk.Bind.bad_of_accept,
    GcZk.Bind.card_coll_le, GcZk.Bind.card_hit₂_le, GcZk.Bind.card_hitOpen_le}
  \leanok
  \uses{def:bind, lem:fresh}
  If the prover is accepted and the extractor does not win, then a fresh first-phase query got the
  answer of an earlier, different one, or a fresh second-phase query got the answer $c$, or the opened point was
  never queried and its value is $c$. Each of these events has probability at most $1/|D|$ for a
  fixed pair, query or point.
\end{lemma}

\begin{proof}
  \leanok
  If the opened point was queried in the first phase, the first query answered with $c$ has a
  different label, and the first occurrence of the opened point is fresh. Otherwise the first
  occurrence of the opened point in the second phase, or the point itself, is fresh. In each case
  apply Lemma~\ref{lem:fresh}.
\end{proof}

\begin{theorem}[Soundness with a hash commitment]
  \label{thm:binding}
  \lean{GcZk.Bind.card_accept_le, GcZk.Bind.prob_accept_le}
  \leanok
  \uses{lem:bad-of-accept}
  The probability that the verifier accepts a cheating prover is at most the probability that the
  extractor outputs the target label, plus
  \[ \frac{q_1^2 + q_2 + 1}{|D|} . \]
\end{theorem}

\begin{proof}
  \leanok
  Sum the bounds of Lemma~\ref{lem:bad-of-accept} over the $q_1^2$ pairs, the $q_2$ queries and the
  opened point.
\end{proof}

\section{End-to-end soundness}

\begin{definition}[Paired prover]
  \label{def:pprover}
  \lean{GcZk.E2E.cans, GcZk.E2E.replayG, GcZk.E2E.replayH, GcZk.E2E.PProver, GcZk.E2E.oH,
    GcZk.E2E.target, GcZk.E2E.RealAccept}
  \leanok
  \uses{def:bind, def:games}
  The cheating prover receives the material and the active labels of its OT input $x$, and at
  each step makes one query to the garbling hash $H$ and one to the commitment oracle $G$. A prover
  that interleaves its queries arbitrarily is of this form after padding. It makes $n_1$ steps,
  commits, sees the opened seed, makes $n_2$ more steps and opens a point. The verifier accepts if
  the point has the target label $L_{\mathrm{out}} + \Delta$ and its value is the commitment.
\end{definition}

\begin{lemma}[Replays]
  \label{lem:replay}
  \lean{GcZk.E2E.replayG_snoc, GcZk.E2E.replayH_snoc, GcZk.E2E.replayG_ans,
    GcZk.E2E.replayH_answers, GcZk.E2E.cans_take, GcZk.E2E.toBind, GcZk.E2E.realAccept_iff}
  \leanok
  \uses{def:pprover}
  For a fixed $H$, the prover answers its own $H$ queries and is a prover of
  Definition~\ref{def:bind} for the oracle $G$, accepted exactly when the paired prover is. For a
  fixed $G$, a strategy that answers its own $G$ queries is an adversary for the oracle $H$.
\end{lemma}

\begin{proof}
  \leanok
  Induction on the number of steps.
\end{proof}

\begin{lemma}[The extractor is an authenticity adversary]
  \label{lem:ext-auth}
  \lean{GcZk.E2E.firstLabel, GcZk.E2E.extAdv, GcZk.E2E.extWins_garbleWin}
  \leanok
  \uses{lem:replay, def:games, thm:eval-garble}
  For a fixed $G$, the extractor, which replays the first phase and outputs the label of the first
  $G$ query answered with the commitment, is an adversary of Theorem~\ref{thm:garble-rom}. If $C$
  is unsatisfiable and the extractor returns the target label, this adversary outputs the inactive
  output label.
\end{lemma}

\begin{proof}
  \leanok
  On an unsatisfiable circuit the active output label is $L_{\mathrm{out}}$, by
  Theorem~\ref{thm:eval-garble}, so the target $L_{\mathrm{out}} + \Delta$ is the inactive output
  label.
\end{proof}

\begin{theorem}[End-to-end soundness]
  \label{thm:e2e}
  \lean{GcZk.E2E.card_accept_le, GcZk.E2E.prob_accept_le}
  \leanok
  \uses{thm:binding, lem:ext-auth, thm:garble-rom}
  Let $|V| \geq 2$, let the seed of the verifier, the hash $H$ and the commitment oracle $G$ be
  uniform, and let $C$ be unsatisfiable. A cheating prover with $n_1$ steps before the commitment and $n_2$ after it is
  accepted with probability at most
  \[ \frac{n_1 + 1}{|V| - 1} + \frac{n_1^2 + n_2 + 1}{|D|} . \]
  In the formalization the hypothesis on $C$ is weaker: $C(x) = 0$ for the input $x$ that the
  prover gives to the OT, which is fixed in advance, with $x_j = 0$ for $j \geq k$.
\end{theorem}

\begin{proof}
  \leanok
  For fixed $H$ and seed, Theorem~\ref{thm:binding} and Lemma~\ref{lem:replay} bound acceptance by
  the extractor's success plus $(n_1^2 + n_2 + 1)/|D|$. For fixed $G$, Lemma~\ref{lem:ext-auth} and
  Theorem~\ref{thm:garble-rom} bound the extractor's success by $(n_1 + 1)/(|V| - 1)$. Sum over all
  the randomness.
\end{proof}
